% sec4_3.tex — Section 4.3: Large-Sample Theory

\section{Large-Sample Theory}
\label{sec:4.3}

Large-sample theory for the single-equation GMM was stated in Propositions~3.1--3.8.
It is just a matter of mechanical substitution to translate it to the multiple-equation
model. Provided that $\boldsymbol{\delta}$, $\boldsymbol{\Sigma}_{\mathbf{xz}}$,
$\mathbf{S}$, $\mathbf{g}_i$, $\mathbf{s}_{\mathbf{xy}}$, and
$\mathbf{S}_{\mathbf{xz}}$ are as defined in the previous sections and that
``Assumption~3.$x$'' is replaced by ``Assumption~4.$x$,''
Propositions~3.1, 3.3, and 3.5--3.8 are valid as stated for the multiple-equation
model. Only a few comments about their interpretation are needed.

\begin{itemize}
\item \textbf{(Hypothesis testing)}\quad
In the present case of multiple equations, $\boldsymbol{\delta}$ is a stacked
vector composed of coefficients from different equations. The import of Propositions~3.3
and~3.8 about hypothesis testing for multiple equations is that we can
test cross-equation restrictions.

\begin{quote}
\textbf{Example 4.2 (continued):}\quad
For the two-equation system of wage equations for two years, the stacked coefficient
vector $\boldsymbol{\delta}$ is
\[
\boldsymbol{\delta} = (\phi_1, \beta_1, \gamma_1, \pi_1,
\phi_2, \beta_2, \gamma_2, \pi_2)'.
\]
It would be interesting to test $\mathrm{H}_0$: $\beta_1 = \beta_2$ and
$\pi_1 = \pi_2$, namely, that the schooling and experience premia remained stable
over time. The hypothesis can be written as $\mathbf{R}\boldsymbol{\delta} = \mathbf{r}$, where
\[
\mathbf{R} = \begin{bmatrix}
0 & 1 & 0 & 0 & 0 & -1 & 0 & 0 \\
0 & 0 & 0 & 1 & 0 & 0 & 0 & -1
\end{bmatrix}, \quad
\mathbf{r} = \begin{bmatrix} 0 \\ 0 \end{bmatrix}.
\]
Nonlinear cross-equation restrictions, too, can be tested; just use part~(c)
of Proposition~3.3 or Proposition~3.8.
\end{quote}

\item \textbf{(Test of overidentifying restrictions)}\quad
The number of orthogonality conditions is $\sum_m K_m$ and the number of coefficients
is $\sum_m L_m$. Accordingly, the degrees of freedom for the $J$ statistic in
Proposition~3.6 are $\sum_m K_m - \sum_m L_m$, and that for the $C$ statistic in
Proposition~3.7 is the total number of suspect instruments from different equations.
\end{itemize}

Proposition~3.2 does not apply to multiple equations as is, but it is obvious how
it can be adapted.

\begin{proposition}[consistent estimation of contemporaneous error cross-equation moments]
\label{prop:4.1}
\textit{Let\/ $\hat{\boldsymbol{\delta}}_m$ be a consistent estimator of\/
$\boldsymbol{\delta}_m$, and let\/
$\hat{\varepsilon}_{im} \equiv y_{im} - \mathbf{z}_{im}' \hat{\boldsymbol{\delta}}_m$ be
the implied residual for $m = 1, 2, \ldots, M$.
Under Assumptions~4.1 and~4.2, plus the assumption that\/
$\E(\mathbf{z}_{im} \mathbf{z}_{ih}')$ exists and is finite for all
$m, h$ $(= 1, 2, \ldots, M)$,}
\begin{equation}
\hat{\sigma}_{mh} \pto \sigma_{mh},
\label{eq:4.3.1}
\end{equation}
\textit{where}
\[
\hat{\sigma}_{mh} \equiv \frac{1}{n} \sum_{i=1}^{n} \hat{\varepsilon}_{im} \hat{\varepsilon}_{ih}
\quad \textit{and} \quad
\sigma_{mh} \equiv \E(\varepsilon_{im}\, \varepsilon_{ih})
\]
\textit{provided that\/ $\E(\varepsilon_{im}\, \varepsilon_{ih})$ exists and is finite.}
\end{proposition}

The proof, very similar to the proof of Proposition~3.2 and tedious, is left as an
analytical exercise.

This leaves Proposition~3.4 about consistently estimating $\mathbf{S}$. The assumption
corresponding to Assumption~3.6 (the finite fourth-moment assumption) is

\begin{assumption}[finite fourth moments]
\label{ass:4.6}
$\E[(x_{imk} \cdot z_{ihj})^2]$ \textit{exists and is finite for
all $k$ $(= 1, 2, \ldots, K_m)$, $j$ $(= 1, 2, \ldots, L_h)$,
$m, h$ $(= 1, 2, \ldots, M)$, where $x_{imk}$ is the $k$-th element of\/
$\mathbf{x}_{im}$ and $z_{ihj}$ is the $j$-th element of\/ $\mathbf{z}_{ih}$.}
\end{assumption}

The multiple-equation version of the formula (3.5.10) for consistently estimating
$\mathbf{S}$ is
\begin{equation}
\widehat{\mathbf{S}} = \begin{bmatrix}
\dfrac{1}{n} \displaystyle\sum_{i=1}^{n}
\hat{\varepsilon}_{i1}\, \hat{\varepsilon}_{i1}\,
\mathbf{x}_{i1} \mathbf{x}_{i1}' & \cdots &
\dfrac{1}{n} \displaystyle\sum_{i=1}^{n}
\hat{\varepsilon}_{i1}\, \hat{\varepsilon}_{iM}\,
\mathbf{x}_{i1} \mathbf{x}_{iM}' \\[8pt]
\vdots & & \vdots \\[4pt]
\dfrac{1}{n} \displaystyle\sum_{i=1}^{n}
\hat{\varepsilon}_{iM}\, \hat{\varepsilon}_{i1}\,
\mathbf{x}_{iM} \mathbf{x}_{i1}' & \cdots &
\dfrac{1}{n} \displaystyle\sum_{i=1}^{n}
\hat{\varepsilon}_{iM}\, \hat{\varepsilon}_{iM}\,
\mathbf{x}_{iM} \mathbf{x}_{iM}'
\end{bmatrix},
\label{eq:4.3.2}
\end{equation}
for some consistent estimate $\hat{\varepsilon}_{im}$ of $\varepsilon_{im}$.

\begin{proposition}[consistent estimation of $\mathbf{S}$, the asymptotic variance of $\bar{\mathbf{g}}$]
\label{prop:4.2}
\textit{Let\/ $\hat{\boldsymbol{\delta}}_m$ be a consistent estimator of\/
$\boldsymbol{\delta}_m$, and let\/
$\hat{\varepsilon}_{im} \equiv y_{im} - \mathbf{z}_{im}' \hat{\boldsymbol{\delta}}_m$ be
the implied residual for $m = 1, 2, \ldots, M$.
Under Assumptions~4.1, 4.2, and~4.6, $\widehat{\mathbf{S}}$ given in
\eqref{eq:4.3.2} is consistent for\/ $\mathbf{S}$.}
\end{proposition}

We will not give a proof; the technique is very similar to the one used in the proof
of Proposition~2.4.

Since this $\widehat{\mathbf{S}}$ is consistent for $\mathbf{S}$, the multiple-equation
GMM estimator, using $\widehat{\mathbf{S}}^{-1}$ as the weighting matrix,
$\hat{\boldsymbol{\delta}}(\widehat{\mathbf{S}}^{-1})$, is an efficient multiple-equation
GMM estimator with minimum asymptotic variance. The formulas for the asymptotic
variance and its consistent estimator are (3.5.13) and (3.5.14), which are reproduced
here for convenience:
\begin{align}
\avar(\hat{\boldsymbol{\delta}}(\widehat{\mathbf{S}}^{-1}))
&= (\boldsymbol{\Sigma}_{\mathbf{xz}}' \mathbf{S}^{-1}
\boldsymbol{\Sigma}_{\mathbf{xz}})^{-1},
\label{eq:4.3.3} \\
\widehat{\avar}(\hat{\boldsymbol{\delta}}(\widehat{\mathbf{S}}^{-1}))
&= (\mathbf{S}_{\mathbf{xz}}' \widehat{\mathbf{S}}^{-1}
\mathbf{S}_{\mathbf{xz}})^{-1}.
\label{eq:4.3.4}
\end{align}

To recapitulate, the features specific to the multiple-equation GMM are:
$\boldsymbol{\Sigma}_{\mathbf{xz}}$ is a block diagonal matrix given by~\eqref{eq:4.1.9},
$\mathbf{S}$ is given by~\eqref{eq:4.1.11},
$\mathbf{s}_{\mathbf{xy}}$ and $\mathbf{S}_{\mathbf{xz}}$ by~\eqref{eq:4.2.2},
and $\widehat{\mathbf{S}}$ by~\eqref{eq:4.3.2}.
For the initial estimator $\hat{\boldsymbol{\delta}}_m$ needed to calculate
$\hat{\varepsilon}_{im}$ and $\widehat{\mathbf{S}}$, we can use the FIVE estimator
(to be presented below) or the efficient single-equation GMM applied to each equation
separately. Of course, as long as it is consistent, the choice of the initial
estimator does not affect the asymptotic distribution of the efficient GMM estimator.

Table~4.1 summarizes the results of this chapter so far.

\begin{table}[htbp]
\centering
\caption{Multiple-Equation GMM in the Single-Equation Format}
\label{tab:4.1}
\small
\medskip
\begin{tabular}{@{}p{3.8cm}p{3.5cm}p{3.5cm}@{}}
\toprule
& Single-Equation GMM applied to the equation in question
& Multiple-Equation GMM \\
\midrule
$\mathbf{g}_i$ & $\mathbf{x}_i \cdot \varepsilon_i$ & (4.1.4) \\[3pt]
$\boldsymbol{\delta}$ & $\boldsymbol{\delta}$ & (4.1.6) \\[3pt]
$\mathbf{s}_{\mathbf{xy}}$
& $\frac{1}{n}\sum \mathbf{x}_i \cdot y_i$ & (4.2.2) \\[3pt]
$\mathbf{S}_{\mathbf{xz}}$
& $\frac{1}{n}\sum \mathbf{x}_i \mathbf{z}_i'$ & (4.2.2) \\[3pt]
Size of $\mathbf{W}$ & $K \times K$
& $\sum_m K_m \times \sum_m K_m$ \\[3pt]
$\boldsymbol{\Sigma}_{\mathbf{xz}}$
& $\E(\mathbf{x}_i \mathbf{z}_i')$ & (4.1.9) \\[3pt]
$\mathbf{S}$ $(\equiv \avar(\bar{\mathbf{g}}))$
& $\E(\varepsilon_i^2\, \mathbf{x}_i \mathbf{x}_i')$ & (4.1.11) \\[3pt]
$\widehat{\mathbf{S}}$
& $\frac{1}{n}\sum \hat{\varepsilon}_i^2\, \mathbf{x}_i \mathbf{x}_i'$
& (4.3.2) \\[3pt]
\midrule
Estimator consistent under which assumptions?
& 3.1--3.4 & 4.1--4.4 \\[3pt]
Estimator asymp.\ normal under which assumptions?
& 3.1--3.5 & 4.1--4.5 \\[3pt]
$\widehat{\mathbf{S}} \pto \mathbf{S}$ under which assumptions?
& 3.1, 3.2, 3.6, $\E(\mathbf{g}_i \mathbf{g}_i')$ finite
& 4.1, 4.2, 4.6, $\E(\mathbf{g}_i \mathbf{g}_i')$ finite \\[3pt]
d.f.\ of $J$ & $K - L$ & $\sum_m(K_m - L_m)$ \\
\bottomrule
\end{tabular}
\end{table}
